% =====================================================================
\section{Introdução}
\label{sec:introducao}

A incorporação de modelos de aprendizado de máquina a sistemas embarcados e a malhas de
controle deslocou uma classe inteira de defeitos para fora do alcance dos métodos
tradicionais de teste. Enquanto um modelo preditivo isolado pode ser avaliado por acurácia
sobre um conjunto de validação, um controlador neural embarcado em uma planta física não
admite esse critério: o que importa não é o erro médio, mas se existe \emph{alguma}
combinação de entradas, dentro do envelope operacional declarado, capaz de levar o sistema
a um estado inseguro. Essa pergunta é de natureza exaustiva, e amostragem --- por mais
densa que seja --- não a responde.

A verificação formal ataca exatamente essa lacuna. No paradigma de \textit{bounded model
checking} (BMC) baseado em SMT, o programa é desenrolado até um limite $K$, traduzido em
uma fórmula lógica e submetido a um \textit{solver} que decide se existe atribuição de
entradas capaz de violar a propriedade especificada. Quando existe, o \textit{solver}
devolve um contraexemplo concreto --- um traço reproduzível que exibe a violação
\cite{cordeiro2012smt}. Aplicada a redes neurais, a técnica vem sendo empregada tanto para
provar robustez local quanto para exibir entradas adversariais concretas
\cite{huang2017robustness}. Quando não existe, obtém-se uma garantia matemática, ainda que
limitada ao horizonte $K$ e às hipóteses declaradas.

Aplicar esse maquinário a componentes de inteligência artificial impõe três dificuldades
específicas. A primeira é de \textbf{escala}: uma camada densa de dimensão $d$ gera
$O(d^2)$ multiplicações simbólicas, e um bloco de \textit{transformer} de produção opera
com $d = 4096$. A segunda é de \textbf{aritmética}: funções de ativação transcendentais
como a sigmoide e a SiLU não possuem axiomatização nas teorias SMT usuais, e o
\textit{solver} as trata como funções não interpretadas, devolvendo \veredito{unknown}. A
terceira, menos discutida na literatura e a que se mostrou mais perigosa neste trabalho, é
de \textbf{especificação}: uma propriedade mal formulada pode ser decidida sem que o modelo
sob verificação participe da prova. O resultado é uma prova verdadeira sobre um objeto que
não é o objeto de interesse --- uma \emph{propriedade vácua}.

Este relatório apresenta o trabalho desenvolvido no ciclo PIBIC 2025--2026 em torno do
verificador ESBMC \cite{gadelha2018esbmc}, um \textit{model checker} de código aberto,
baseado em SMT, com suporte a C, C++ e ponto flutuante. O projeto percorreu sete estudos de
caso de complexidade crescente, do neurônio quantizado isolado à malha fechada completa
formada por planta e controlador de aprendizado por reforço, e produziu, além dos vereditos,
uma infraestrutura de reprodutibilidade: um invólucro Python sobre o binário do verificador,
uma suíte de regressão que impede a confusão entre erro de execução e propriedade violada, e
um gerador de evidências que reexecuta cada \textit{harness} citado no texto e versiona a
saída bruta do \textit{solver}.

A justificativa para essa ênfase em procedência de evidência não é retórica. Uma auditoria
sistemática do próprio repositório, conduzida na fase final do ciclo, encontrou afirmações
quantitativas que os artefatos não sustentavam: tempos de verificação que provinham de
valores fixos em \textit{scripts} de plotagem, um experimento de escalabilidade cujo
parâmetro de dimensão nunca chegava ao código compilado, um orquestrador de verificação com
o portão de chamada ao \textit{solver} comentado, e duas propriedades de segurança que
decidiam sem tocar nos pesos da rede. Nenhum desses defeitos produzia erro visível; todos
produziam números plausíveis. A correção deles, documentada ao longo da
Seção~\ref{sec:resultados}, é parte substantiva do resultado deste PIBIC.

O relatório está organizado da seguinte forma. A Seção~\ref{sec:objetivos} recupera os
objetivos do plano de trabalho e indica o grau de cumprimento de cada um. A
Seção~\ref{sec:fundamentacao} apresenta os fundamentos de \textit{model checking} limitado e
indutivo, de codificação SMT e de quantização em ponto fixo. A Seção~\ref{sec:metodologia}
descreve a infraestrutura construída, o protocolo de evidência e a formulação de cada
estudo de caso. A Seção~\ref{sec:resultados} reporta e discute os resultados medidos. A
Seção~\ref{sec:limitacoes} explicita limitações, ameaças à validade e dificuldades
enfrentadas. A Seção~\ref{sec:conclusao} sintetiza as conclusões e a
Seção~\ref{sec:cronograma} registra o cronograma efetivamente executado.
